\chapter{Occupational Therapy}

\section*{Overview}
Occupational therapy is a rehabilitation service that helps a person perform everyday activities safely and independently after illness, injury, disability, or developmental difficulty.

\section*{Why It Is Done}
It may help with arthritis, stroke, neurological conditions, hand injuries, chronic pain, developmental conditions, or recovery after hospitalization. The therapist focuses on activities meaningful to the person's home, school, work, and community life.

\section*{How It Is Performed}
The therapist evaluates movement, strength, coordination, cognition, sensation, and daily tasks. Treatment may include guided exercises, practice of self-care, energy-conservation strategies, adaptive equipment, splinting, and changes to the home or work environment.

\section*{Preparation and Safety}
Wear comfortable clothing and bring information about symptoms, medicines, and activities that are difficult. Exercises should be progressed gradually. Tell the therapist about pain, dizziness, new weakness, or worsening symptoms.

\section*{Results and Follow-up}
Progress is measured by improvements in function, safety, range of motion, strength, and ability to complete daily activities. The therapist may provide a home program and coordinate care with physicians, physical therapists, speech therapists, and caregivers.

\section*{References}
\begin{enumerate}
  \item MedlinePlus. Occupational therapy. U.S. National Library of Medicine; 2025.
  \item Current Medical Diagnosis and Treatment. McGraw-Hill; 2025.
\end{enumerate}
\clearpage
